\chapter{Lezione 11/03}

In una struttura di Kripke, la funzione L associa a ogni stato un modello proposizionale che rappresenta le proprietà atomiche vere in quello stato.
Possiamo quindi definire un concetto di soddisfacibilità locale, ovvero in uno stato.


Dobbiamo comunque avere dei sistemi per mappare i sistemi in grafi. Per sistemi complessi sarà necessario decomporre il sistema in buildibg block più semplici.
I building block \q{atomici} sono i Program Graph, che sono definiti da un insieme di variabili \(\Var = \set{v_1,\ldots,v_k}\) con dominio \(D = \bigcup_{1,\ldots,k} D_i\), con \(D_i\) dominio di \(v_i\).

In particolare un Program Graph è una sestupla \((\Loc,\Act,\Eff,\hookrightarrow , \Loc_0, g_0)\) dove:
\begin{itemize}
  \item \(\Loc\) è un insieme di locazioni o stati di controllo, (dominio del program counter)
  \item \(\Act\) l'insieme delle azioni che modificano lo stato del programma
  \item \(\Eff : \Act\times \Eval(\Var) \to \Eval(\Var)\) è una funzione
  \item \(\hookrightarrow \subseteq \Loc \times \Cond(\Var) \times \Act \times \Loc\) è la relazione di transizione
  \item \(\Loc_0 \subseteq \Loc\) è l'insieme delle locazioni iniziali
  \item \(g_0 \in \Eval(\Var)\) è una condizione che le variaibili devono soddisfare all'inizio dell'esecuzione
\end{itemize}
La relazione di transizione \(\hookrightarrow\) è una relazione che collega due locazioni tramite una condizione sulle variabili, (in un certo senso collegando \q{temporalmente} locazioni o righe del programma),detta anche guardia ovvero la verifica di una certa proprietà sulle variabili, e un'azione.
La guardia mi dice se è possibile eseguire la transizione, e l'azione mi dice come modificare le variabili quando eseguo la transizione.

La funzione \(\Eff\) invece descrive la semantica delle azioni: prende in input un'azione e deve associare una trasformazione del valore delle variabili prima e dopo l'azione. In altre parole collega \q{temporalmente} stati delle variabili.
Per fare ciò è necessario introdurre un concetto di valutazione delle variabili, e per farlo usiamo \(\Eval\), un insieme di funzioni che mappano ogni variabile al suo valore, rendendo \(\Eval\) l'insieme di tutte le possibili configurazioni delle variabili.

\(\Act\) e \(\Cond\) sono definite sulla base del concetto di espressioni \(\Expr(\Var\cup D)\). Questa dipende dai domini delle variabili, se prendiamo ad esempio le espressioni aritmetiche abbiamo \(v+1\).
Le condizioni sono dunque valutazioni tra espressioni del tipo \(\expr_1 \cdot \expr_2\), con \(\cdot\) dipendende dal dominio (per aritmetica \(=, \neq, <, > ecc\)), eventualmente anche con operatori logici.
Le azioni su \(\Var\) è l'insieme degli assegnamenti, ovvero espressioni del tipo \(v := \expr\).

\todo{Esempio vending machine.}

Possiamo notare che i grafi sottostanti a questo e all'esempio precedente sono isomorfi, ma i due ogetti sono diversissimi, poiché il secondo sta descrivendo un LTS con stati infiniti (essendo il dominio delle variabili infinito).
È possibile però \q{compilare} un ogetto del genere in una Kripke Structure, ma è necessaria la funzione \(\Eff\) per definire la semantica delle azioni.
Questa compilazione di fatto \q{da semantica} a questa struttura.

Per lo spazio degli stati questo equivale ad \(\Eval(\Var)\) (variabili esplicite) e \(\Loc\) (variabili implicite, ovvero il program counter).

Formalmente una valutazione in \(\Eval\) è una funzione \(\eta: \Var \to D\), ognuna di questa può essere rappresentata con una tupla lunga col numero di variabile, che accoppiata con una locazione \(l\) rappresenta uno stato del sistema, e dunque un nodo del grafo orientato sottostante.

Quindi nella Kripke Structure sottostante, \(S = \Loc \times \Eval(\Var)\).

Col concentto di \(\Eval(\Var)\) definito possiamo parlare meglio di \(\Eff\). Data \(\eta \in \Eval(\Var)\) e un azione della forma \(a \overset{\text{def}}{=} v := \expr\), con \(\expr \in \Var \cup D\) allora l'effetto dell'azione è
\[\Eff(a,\eta)=\eta[v\leftarrow \expr]\]
con
\[\eta[v\leftarrow \expr] = \eta'\]
dove
\[\eta' (x)= \begin{cases}
    \expr   & \text{se } x = v  \\
    \eta(x) & \text{altrimenti}
  \end{cases}\]

Questa definizione non è ben definita perché ci vuole un concectto di valutazione dell'espressione \(\expr\) in \(\Eval(\Var)\). In altre parole un \(\Val(\expr, \eta) \in D\), da definire per induzione strutturale sulla struttura delle espressioni.
I casi base sono costanti e le variabili (valutate con \(\eta\)). Il resto dipende dalla semantica degli operatori.

Quindi per essere precisi la definizione di \(\Eff\) è la seguente:
\[\Eff(a,\eta)=\eta[v\leftarrow \Val(\expr, \eta)]\]
con
\[\eta[v\leftarrow \Val(\expr,\eta)] = \eta'\]
dove
\[\eta' (x)= \begin{cases}
    \Val(\expr,\eta) & \text{se } x = v  \\
    \eta(x)          & \text{altrimenti}
  \end{cases}\]

Per definire ora la funzione di transizione dovremmo definire quali stati sono collegati da una transizione basandoci su \(\hookrightarrow\) e \(\Eff\).

In particolare definiamo la relazione di transizione \(\xrightarrow{} \subseteq S \times \Act \times S\) tale che:
\begin{itemize}
  \item Se \(l \xhookrightarrow{g:a} l' \) è nel \(PG\) e \(\eta \models g\) allora \(\angle{l,\eta} \xrightarrow{a} \angle{l',\eta'}\) è nella relazione di transizione, con \(\eta' = \Eff(a,\eta)\).
\end{itemize}

Gli stati iniziali sono definiti come \(S_0 = \set{\angle{l, \eta}[l \in \Loc_0 \land \eta \models g_0]}\)

Infine per l'etichettatura, definiamo con \(AP = \Loc \cup \Cond(\Var)\) e \(L : S \to 2^{AP}\) tale che \(L(\angle{l,\eta}) = \set{l} \cup \set{g \in \Cond(\Var)}{\eta \models g}\).

Fatto ciò, abbiamo definito il linguaggio per definire i componenti, ma questo è poco più di un programma sequenziale. Un po' di più perché da possibilità di scelte non deterministiche.
Per sistemi complessi è però molto difficile specificarli per intero.

Dunque l'approccio più comunque è quello di specificali tramite composizione di componenti più semplici, che interagiscono in modo sequenziale o concorrente.
La concorrenza è modellabile in tantissimi modi diversi. Esempi da trascrizione.

La composizione parallela sincrona è facile da definire. Dati \(TS_1, \ldots, TS_n\) sistemi a transizione tali che \(TS_i = (S_i, S_{0i}, A_i, R_i)\) con \(S_i\) disgiunti, allora la composizione parallela sincrona è il sistema a transizioni \(TS = (S,S_0,R,L)\) con:
\begin{itemize}
  \item \(S = S_1 \times \ldots \times S_n\)
  \item \(S_0 = S_{01} \times \ldots \times S_{0n}\)
  \item \(A = A_1 \cup \ldots \cup A_n\)
  \item \(R \subseteq S \times A \times S\) è taòe che, dati \(\angle{s_1, \ldots, s_n}, \angle{s_1', \ldots, s_n'} \in S\) e \(\angle{a_1,\ldots, a_n} \in A\) contiene \(\angle{s_1,\ldots,s_n}\xrightarrow{\angle{a_1,\ldots,a_n}}\angle{s_1',\ldots,s_n'}\) se \(s_i \xrightarrow{a_i} s_i', \forall i \in \set{1,\ldots,n}\).
\end{itemize}

\chapter{Lezione 16/03}

Il problema delle espressioni regolari è che non hanno esplicitamente una congiunzione o intersezione.
Se avessi l'intersezione potrei scrivere \(\bigwedge_{i=0}^n F p_i\) come \(\bigcap_{i=1}^n \Sigma* p_i \Sigma^*\).

Non avendolo però è necessario scrivere un'espressione regolare per ogni ordine.

Questo è perché nei linguaggi regolari non c'è il complemento, e dunque non c'è l'intersezione, questo perché di base non è necessario.
Ciò comporta però che descrivere una proprietà che non è presente sia complicato.

Quindi vogliamo un linguaggio che si avvicini di più all'idea informale di proprietà, poiché è già difficile e inevitabilmente inesatto il processo di formalizzazione di un idea.
Se il linguaggio è più vicino all'idea informale, allora sarà più facile formalizzare correttamente l'idea informale, e dunque sarà più facile verificare che il sistema soddisfa la proprietà desiderata.

Passiamo quindi a un linguaggio logico. Inizialmente con il prim'ordine, per poi passare alla logica temporale, che cerca di nascondere all'interno di operatori logici dei pattern ricorrenti in logica del prim'ordine.

Quali sono gli ogetti che ci interessano? L'universo, o ogetti del discorso, sono le parole \(w: [0,n)\to \Sigma\) oppure \(w: \N \to \Sigma\).

Immaginiamo di voler esprimere la proprietà \qi{Prima o poi il sistema ha stampato ed è in attesa del nuovo input}.
Essendo una congiunzione di proprietà, il che però deve far si che gli elementi dell'alfabeto non sia da considerarsi come atomiche.
Dobbiamo considerare \(\Sigma = 2^{AP}\), con \(AP\) l'insieme delle proposizioni atomiche.

\todo{Espandi esepio con \(p\) e \(i\) print e idle}

Dunque i nostri modelli saranno essenzialmente \(w:\N \to 2^{AP}\), ovvero a ogni istante temporale associamo l'insieme di proposizioni atomiche vere in quel momento.

Vediamo dunque col prim'ordine possiamo esprimere queste proprietà.
Dobbiamo quindi avere \(w \models \phi\).

Per il prim'ordine avremo una grammatica di questo tipo, con \(p \in AP\):
\[\phi ::= p(x) \mid \neg \phi \mid \phi \land \phi \mid \phi \lor \phi\mid \forall x \phi \mid \exists x \phi\]

Normalmente una grammatica di questo tipo porterebbe a un problema indecidibile, ma se ci concentriamo su sole stringhe come modelli allora torna decidibile.

Vediamo dunque la semantica del prim'ordine. La definiremo in maniera composizionale, ovvero definendo la semantica di una proposizione atomica e come gli operatori logici modificano la semantica delle proposizioni atomiche.
In particolare ci serve un assegnamento \(\chi : \Var \to N\) che ci dice ogni variabile a quale istante temporale si riferisce.
\(w,\chi \models p(x)\) se \(p \in w(\chi(x))\). Con \(w(\chi(x))\) si intende l'elemento di \(w\) in posizione \(\chi(x)\), ovvero l'insieme di proposizioni atomiche vere in quel momento, e \(p \in w(\chi(x))\) significa che \(p\) è una di queste proposizioni atomiche vere in quel momento.

Prendiamo ad esempio la parola \(w = \emptyset \set{i} \set{i} \set{i} \emptyset \set{i,p}\). Allora con \(\chi = \set{x \to 3}\) avremo \(w,\chi \not\models p(x)\) poiché \(p \not\in w(3)\), mentre con \(\chi' = \set{x \to 5}\) si ha \(w,\chi' \models p(x)\) poiché \(p \in w(5)\).

Continuando la semantica, avremo che:
\(w,\chi \models \neg \phi\) se\(w,\chi \not\models \phi\).
\(w,\chi \models \phi_1 \land \phi_2\) se \(w,\chi \models \phi_1\) e \(w,\chi \models \phi_2\).
\(w,\chi \models \phi_1 \lor \phi_2\) se \(w,\chi \models \phi_1\) o \(w,\chi \models \phi_2\).
Questo chiude la semantica della parte proposizionale.

Ora per dare la semantica agli operatori quantificatori, diremo:
\(w,\chi \models \forall x.\phi\) se per ogni \(n \in \N\), \(w,\chi[x \to n] \models \phi\). Ovvero sostituendo a \(x\) ogni possibile istante temporale, la formula \(\phi\) è soddisfatta.
\(w,\chi \models \exists x.\phi\) se esiste \(n \in \N\) tale che \(w,\chi[x \to n] \models \phi\). Ovvero sostituendo a \(x\) almeno un istante temporale, la formula \(\phi\) è soddisfatta.

Vediamo come modellare una proprietà sul futuro. Tipo \(\exists x. i(x) \land p_2(x)\). Ad esempio questo può essere inteso come \qi{Prima o poi il sistema ha stampato ed è in attesa del nuovo input}.
Se invece voglio dire che non ha mai due processi in una sezione critica, allora posso scrivere \(\forall x. \neg (p_1(x) \land p_2(x))\).

Se invece voglio dire che infinitamente spesso valga la proprietà \(i\). Qui ho una difficoltà, perché mi serve poter esprimere una relazioni tra istanti temporali.
Ci serve quindi poter esprimere il concetto di \qi{dopo}, è necessario introdurre una relazione d'ordine tra istanti temporali.
Sintatticamente inroduciamo il simbolo \(\preceq\), la cui semantica è:
\(w,\chi \models x\preceq y\) se \(\chi(x) \leq \chi(y)\). In questo caso sono stati nettamente separati i simboli \(\preceq\) e \(\leq\) per evidenziare che il primo è un simbolo del linguaggio ogetto, mentre l'altro è l'ordine usuale dei numeri naturali.
Nella pratica useremo \(\leq\) anche nel linguaggio ogetto poiché sarà chiaro dal contesto se si tratta di un simbolo del linguaggio o dell'ordine usuale dei numeri naturali.

Per esprimere quindi che una proprietà \(i\) si verifica infinitamente spesso, possiamo scrivere \(\forall x. \exists y. x \leq y \land i(y)\).
Se voglio esprimere che arriva un punto dopo il quale sarò sempre in idle posso scrivere \(\exists x. \forall y. x \leq y \to i(y)\).

In generale però possiamo osservare che quando parliamo di proprietà temporali ci troviamo sempre in uno di questi due casi:
\begin{itemize}
  \item \qi{Prima o poi} \(p\) è vero, ovvero \(\exists y. x \leq y \land p(y)\)
  \item \qi{Dopo un certo punto} \(p\) è sempre vero, ovvero \(\exists x. \forall y. x \leq y \to p(y)\)
\end{itemize}

Ma allora possiamo costruire un linguaggio che abbia sintatticamente due operatori che esprimono esattamente questi due pattern ricorrenti come semantica.

\section{Linear Temporal Logic}

Introdotta da Pnueli nel 1979. In questa logica lasciamo implicito l'istante in cui si verifica una proprietà, ci penserà la semantica degli operatori a gestire questo aspetto.
La grammatica è la seguente, con \(p \in AP\):
\[\phi ::= p \mid \neg \phi \mid \phi \land \phi \mid \phi \lor \phi \mid F\phi \mid G\phi\]
Dove \(F\) è l'operatore \q{prima o poi} (o \q{eventually}, \q{future}) e \(G\) è l'operatore \q{dopo un certo punto} (o \q{always}, \q{globally}).

Ovviamente anche ora parleremo di proprietà della parole, ma ora a differenza del prim'ordine non useremo più un assegnamento. Invece, per la semantica, utilizzeremo un indice \(i\) che rappresenta l'istante temporale in cui stiamo valutando la formula, e di conseguenza da cui iniziamo a guardare nel futuro.

La semantica è la seguente:
\(w,i \models p\) se \(p \in w(i)\).
La semantica degli operatori proposizionali è identica.
Venendo ai nuovi operatori invece abbiamo:
\(w,i\models F\phi\) se \(\exists j \geq i \st w, j \models \phi\), ovvero se \(\exists j: j \geq i \land w,j\models\phi\).
\(w,i\models G\phi\) se \(\forall j \geq i , w, j \models \phi\), ovvero se \(\forall j: j \geq i \to w,j\models\phi\).
Con questa sintassi però non è possibile esprimere proprietà che mi descrivono \q{il passato}, come \(\exists x \st p_1(x) \land \forall y, y\leq \neg p_2(x)\).

Di conseguenza, alla sintassi nell'81 vennero sostituiti \(F\) e \(G\) con due operatori binari e dimostrato che con questa agginuta la logica temporale risulta essere equivalente al prim'ordine.
In particolare, la sintassi diventa:
\[\phi ::= p \mid \neg \phi \mid \phi \land \phi \mid \phi \lor \phi \mid \phi U \phi \mid R\phi\]
Dove \(U\) è l'operatore \q{until} e \(R\) è l'operatore \q{release}
Intuitivamente \(U\) esprime che una proprietà \(p\) è vera fino a quando non diventa vera una proprietà \(q\), mentre \(R\) esprime che una proprietà \(p\) è vera a meno che non diventa vera una proprietà \(q\).

Formalmente:
\(w,i \models \phi_1 U \phi_2\) se \(\exists j \geq i \st w,j\models \phi_2\) e \(\forall k: i \leq k < j \to w,k\models \phi_1\)
\(w,i \models R\phi_1 \phi_2\) se \(\forall j \geq i \st w,j\models \phi_2\) o \(\exists k: i \leq k < j \land w,k\models \phi_1\)

Possiamo notare che \(F\phi \equiv \top U \phi\) e che \(G\phi \equiv \bot R \phi\), quindi \(F\) e \(G\) sono esprimibili in termini di \(U\) e \(R\).

\chapter{Lezione 18/03}

Se da un lato ci sono i modelli completamenti sincroni dall'altra ci sono quelli completamente asincroni, che generalizzano quelli in cui alcuni si sincronizzano e altri no.

Per farlo in maniera semplice possiamo usare il modello di composizione sincrona ma aggiungendo ad ogni insieme di azioni locali una no operation \(-\) per permettere a ogni automa locale di effettuare un'azione vacua.

Va ovviamente rilassata nella definizone della relazione di transizione globale il vincolo che ogni azione globale sia composta da azioni locali, e invece dire che ogni azione globale è composta da azioni locali o no operation, ma in questo caso lo stato locale di partenza e lo stato locale di arrivo devono essere uguali.

Questo porta chiaramente ad essere raggiungibili stati che precedentemente poteva non esserlo. Questo modello però permette di esprimere sia comportamenti sincroni che asincroni che comportamenti intermedi, e dunque è più generale.

Inoltre questo tipo di modello con asincronismo virtuale è esattamente come vengono gestiti nella pratica i processi concorrenti, infatti come sott'insieme di questo modello abbiamo la composizione asincrona con interleaving, in cui c'è composizione completamente asincrona e ogni azione globale è composta da una sola azione locale e no operation per gli altri stati in maniera non deterministca.

Ovviamente in questo modello non è inficiata la raggiungibilità di stati, ma non vi è vero parallelismo.

\todo{Esempio <ok,ok>, }

Una formulazione equivalente del modello asincrono con interleaving più compatto è quella di avere come insieme delle azioni l'unione disgiunta degli insiemi di azioni locali, ad esempio avendo le azioni rappresentate come coppie \(\angle{i,a}\) con \(i\) indice dell'automa locale e \(a\) azione locale. Allora possiamo appropriamente riformulare la relazione di transizione come da slide.

Questo è un introduzione al concetto di composizione sincrona e asincrona (con interleaving).

Noi avevamo però dei Program Graph per descrivere i componenti, che poi venivano \q{compilati} in LTS. Abbiamo visto come fare composizione di LTS, ma per i program graph dobbiamo modellare la possibilità di poter condividere informazioni.

Nel mondo degli LTS però non c'è il concetto di variabile. Quindi l'approccio di tradurre vari PG in LTS separatamente per poi comporli non è necessariamente corretto, poiché nel passaggio da PG a LTS si perde il concetto di variabile e dunque di condivisione di variaibli.

Dunque questa fusione va fatta a livello di PG, ma allora la composizione va fatta a livello di PG, e dunque è necessario definire una composizione di PG.

Interleaving Program Graph

Dato un insieme di program graph, definiamo l'insieme delle variabili globali come l'unione di tutte le variabili dei PG (NON DISGIUNTA)\footnote{Si assume che se due variabili hanno lo stesso nome sono condivise, se non lo erano andavano disambiguate prima. Una definizione alternativa tiene conto di variabili non condivise con lo stesso nome ma è inutilmente più complesso}.

L'insieme delle locazioni, come le locazioni iniziali sono il prodotto cartesiano. Per le azioni possiamo tranquillamente avere l'unione disgiunta. Per le guardie iniziali invece abbiamo la congiunzione di tutte le guardie iniziali.

Dobbiamo dunque definire le transizioni, ma sono definite come negli LTS, mentre per \(\Eff\) è più complesso.

Abbiamo a disposizione solo le \(\Eff_i: \Act_i \times \Eval(\Var_i)\to \Eval(\Var_i)\). Dobbiamo definire \(\Eff_i:\biguplus \Act_i \times \Eval(\bigcup\Var_i)\to \Eval(\bigcup\Var_i) \).

Andiamo a definire le restrizioni delle \(\eta\), \(\eta_i\).

Troppo veloce, guarda slide.

Esempi mutua esclusione.

C'è un problema di fairness. A noi non interessa come uno scheduler implementi la fairness, purché mi garantisca una certa fairness.

Dobbiamo capire come introdurre dei vincoli di fairness, senza però inserire nel modello uno specifico scheduler.

Per fare ciò quindi non possiamo andare a modificare l'LTS, piuttosto dobbiamo avere un filtro sulle possibili esecuzioni, avendo quindi un algoritmo di analisi che ignori le computazioni che non rispettano i vincoli di fairness.

Ci sono tipi diversi vincoli di fairness:
\begin{description}
  \item[Fairness incondizionata:] ogni processo dev'essere schedulato infinitamente spesso.
  \item[Fairness debole:] se un processo è continuamente abilitato, allora deve essere schedulato infinitamente spesso.
  \item[Fairness forte:] se un processo è abilitato infinitamente spesso, allora deve essere schedulato infinitamente spesso.
\end{description}

La fairness incondizionata è troppo debole, poiché non garantisce che quando un processo viene schedulato sia abilitato a fare qualcosa. La fairness debole è più ragionevole, ma ancora non è sufficiente, poiché richiede che al processo è richiesto di essere continuamente abilitato, che è una richiesta molto forte.

\chapter{Lezione 25/03}

Un concetto a cui bisogna stare attenti durante la modellazione è quello dell'atomicità. Un'azione atomica si traduce banalmente in una transizione, ma ciò implica che si è modellato adeguatamente in modo che il sistema concreto possa trattare come atomico ogni cosa che è stato modellato con un unica transizione.

Un tipico esempio è l'operazione di incremento, che potrebbe essere atomica come no, quindi in base al sistema concreto è possibile modellarlo con una transizione o meno.

Approfondiremo più avanti.

Esistono anche contesti di comunicazione con scambio di messaggi. Se stiamo modellando una singola macchina, questa è effettivamente implementata usando una stessa memoria, e quindi le variabili condivise sono equivalenti, ma se stiamo parlando di diverse macchine già è diverso.

L'idea per modellare questo è usare dei simboli condivisi di comunicazione che etichettano le azioni che rappresentano l'invio o la ricezione di un messaggio.

Quindi definisco \(Sync = A_1 \cap A_2\) insieme delle azioni sincronizzanti. Affinché uno dei due possa fare un azione in \(Sync\) c'è bisogno che entrambi la facciano, in modo che quell'azione sia in parallelismo reale.

Tale modello è detto di sincronizzazione via handshake.

La differenza sul \(TS\) composto è che l'insieme di azioni \(A \subset Sync \cup (A\cup\set{-}\setminus Sync)\times(A\cup\set{-}\setminus Sync)\) e la relazione di transizione è tale che \(\langle s_1,s_2\rangle \xrightarrow{a} \langle s_1',s_2'\rangle\) se \(a \in Sync\) e \(s_1 \xrightarrow{a} s_1'\) e \(s_2 \xrightarrow{a} s_2'\) oppure \(a = \langle a_1, a_2\rangle \in A\) e \(s_o\xrightarrow{a_1} s_i'\) per \(i \in \set{1,2}\) se \(a_i \neq -\) o % completa

Un modo tipico è avere handshake con arbitro, in cui compongo oltre ai due processi anche un arbitro che gestisce la comunicazione, e dunque i processi non comunicano direttamente tra loro ma tramite l'arbitro.

Per portarlo a livello di PG si parla di Channel system. Sono essenzialmente dei program graph, in cui le variabili sono tutte locali ma ho dei canali di comunicazione tra i processi.

\chapter{Lezione 30/03}

LTL ci permette di esprimere praticamente tutto ciò che è esprimibile con le espressioni \(\omega\)-regolari, tranne cose come il counting.

Dire che una proprietà \(P\) è vera in posizioni pari (NON ESCLUSIVAMENTE SUI PARI) si può fare facilmente con \({((\Sigma\Sigma)*(p\Sigma))}^\omega\) ma non è possibile esprimere questa proprietà in LTL, come neanche in prim'ordine.

Comunque in entrambe le sintassi manca anche l'operatore di next, che però ha senso solo se si parla di modelli con tempo discreto, e con modelli con tempo continuo non è detto che abbia senso.

Tale operatore è espresso come \(X \phi\), e la sua semantica è \(w,i \models X\phi\) se \(w,i+1 \models \phi\).

Grazie a questo operatore possiamo dire cose come \(G(p\land X \neg p)\), che può sembrare descriva la proprietà di essere vero solo in posizioni pari.

Ma andando a sviluppare questa formula, otteniamo:
\[w,0\models G(p\land X \neg p) \iff \forall j \geq 0, w,j \models (p\land X \neg p) \iff \forall j \geq 0, (w,j \models p \land w,j+1 \models \neg p)\]

Che è chiaramente insoddisfacibile. Il problema è che questa formula pretende che \(p\) sia vero in ogni punto.
Un modo più corretto per esprimere questa proprietà è \(G((p \leftrightarrow X \neg p) )\). Ad essere precisi questa ci dice che \(p\) è vera o solo sui pari o solo sui dispari. Per ottenere la proprietà che \(p\) è vera sui pari e falsa sui dispari, allora basta innescarla con \(p\) all'inizio, ovvero \(p \land G((p \leftrightarrow X \neg p) )\).

Espandendola abbiamo:
\[w,0 \models p \land G((p \leftrightarrow X \neg p) )\iff w,0 \models p \text{ e } \forall j \geq 0, w,j \models p \iff w,j+1 \models X \neg p\]

In questo modo abbiamo che \(p\) è vera in \(0\), allora deve essere falsa in \(1\), allora deve essere vera in \(2\), e così via, ottenendo la proprietà desiderata.

Questa però ci dice che la proprietà è soddisfatta solo in posizioni pari, non che è soddisfatta almeno in posizione pari, che abbiamo già detto non è esprimibile in LTL\@.

Dobbiamo ora dare una nozione formale al meta livello di equivalenza tra formule.

Ad esempio abbiamo che \(\land\) e \(\lor\) rispettano le leggi di De Morgan.

Inoltre abbiamo che \(U\) e \(R\) sono duali, infatti \(w,i \models \phi_1 U \phi_2\) se e solo se \(w,i \models \neg (\neg \phi_1 R \neg \phi_2)\).

Andiamo a dimostrare questa equivalenza:

\[w,i\models \phi_1 U \phi_2 \iff \exists j \geq i( w,j \models \phi_2 \land \forall i\leq k < j, w,k\models \phi_1)\iff\]
\[\iff \neg \neg \exists j \geq i\st( w,j \models \phi_2 \land \forall i\leq k < j, w,k\models \phi_1) \iff\]
\[\iff \neg \forall j \geq i,\neg( w,j \models \phi_2 \land \forall i\leq k < j, w,k\models \phi_1) \iff\]
\[\iff \neg \forall j \geq i, w,j \models \neg \phi_2 \lor \neg \forall i\leq k < j, w,k\models \phi_1 \iff\]
\[\iff \neg \forall j \geq i, w,j \models \neg \phi_2 \lor \exists i \leq k < j\st w,k\models \neg \phi_1 \iff\]
\[\iff w,i \models \neg (\neg \phi_1 R \neg \phi_2)\]

Ovviamente per essere un equivalenza stiamo dicendo che tutti questi i passaggi siano corretti a prescindere da \(w\) e \(i\).
In questo caso in effetti non abbiamo usato nulla di specifico su \(w\) e \(i\), e quindi indipendentemente dalla scelta dalla parola e dal punto in cui partiamo, questa equivalenza è sempre vera.

\begin{definition}{Equivalenza tra formule}
  Diremo che due formule \(\phi\) e \(\psi\) sono equivalenti, e scriveremo \(\phi \equiv \psi\), se per ogni parola \(w\) e ogni punto \(i\) si ha che \(w,i \models \phi\) se e solo se \(w,i \models \psi\).
\end{definition}

Grazie a questa definizione possiamo mostrare come LTL del 71 sia sussunta da quella dell'81, infatti \(F\phi \equiv \top U \phi\) e \(G\phi \equiv \bot R \phi\), dove \(\top\) e \(\bot\) sono simboli sintattici assunti presenti nella grammatica, con semantica \(w,i \models \top\) sempre vero e \(w,i \models \bot\) sempre falso. Questo era facile da intuire, poiché la semantica di \(F\) e \(G\) è parte di quella di \(U\) e \(R\), che non fanno altro che aggiungere ulteriori vincoli, i quali vengono ignorati se si usano \(\top\) e \(\bot\). Ovviamente questo mostra una proprietà di dualità tra \(F\) e \(G\), che è la stessa dualità tra \(U\) e \(R\).

Molto spesso, comunque, quello che mi interessa non è usare il counting, quindi il fatto che LTL non sia in grado di esprimere proprietà di counting non è un grosso problema. Inoltre, l'operatore \(X\) non è qualcosa di assoluto, poiché dipende dalla granularità della discretizzazione del tempo. Molto spesso a livello di proprietà pratiche non si usa proprio.

Omettendolo ottengo il frammento stattering free di LTL, che non è espressivamente equivalente a LTL intero, poichè non è possibile distinguere tra parole in cui c'è \qi{idle} in una posizione. Ad esempio le due parole \(ppp\neg p \neg p pp\) e \(p\neg p p\) sono indistinguibili in questo frammento, ma distinguibili in LTL\@. In un certo senso non importa \qi{quanto tempo} sto in uno stato, ma solo i cambiamenti di stato. La cosa importante di questa cosa è che questo operatore porta quasi sempre ad un aumento della complessità da \(NP\) a \(PSPACE\), quindi è un operatore che è meglio evitare se non è necessario.

Proprietà di one-step unfolding. Si può facilmente vedere che vale \(F\phi\equiv \phi \lor X F \phi\). Ho preso la mia formula e l'ho \qi{srotolata} di un passo. In teoria posso srotolarla di un numero qualsiasi di passi, ma srotolarla di un passo è sufficiente per poter costruire un automa. Vale lo stesso per \(G\phi\equiv\phi\land XG\phi\). Dunque saranno vere le opportune generalizzazioni per \(U\) e \(R\). In particolare

\[\phi_1 U \phi_2 \equiv \phi_2 \lor (\phi_1 \land X(\phi_1 U \phi_2 ))\]
\[\phi_1 R \phi_2 \equiv \phi_2 \land (\phi_1 \lor X(\phi_1 R \phi_2 ))\]

NNF ugale a logica classica usando dualità tra \(U\) e \(R\) e tra \(F\) e \(G\). Per quanto riguarda \(X\) invece, è un operatore self-dual, ovvero \(\neg X \phi \equiv X \neg \phi\).

Scrivi algoritmo formale (banale).

\begin{definition}{Lunghezza}
  la lunghezza di una formula \(\phi\) è una funzione che rispetta le seguenti regole:
  \begin{itemize}
    \item \(len(p) = 1\) per \(p \in AP\)
    \item \(len(Op \phi) = 1 + len(\phi)\) per \(Op \in \set{\neg,X}\)
    \item \(len(\phi_1 Op \phi_2) = 1 + len(\phi_1) + len(\phi_2)\) per \(Op \in \set{\land, \lor, U, R}\)
  \end{itemize}
\end{definition}

Chiaramente alla peggio portare in NNF mi raddoppia la lunghezza della formula, quindi se la formula originale è di lunghezza \(n\) allora la formula in NNF è di lunghezza al più \(2n\). Questo perché il peggior caso e qullo in cui abbiamo \(\neg\) come opertaore più esterno e tutti i letterali positivi. In questo caso infatti considerando l'albero sintattico andiamo a togliere il nodo in radice e aggiungere un nodo per ogni foglia. Essendo che il numero di foglie è lineare rispetto alla lunghezza della formula, allora la lunghezza della formula in NNF è al più \(2n\).

Vediamo ora qualche esempio di proprietà espressa di modellazione di sistema espressa in LTL\@.

\paragraph{Reachability}

Classicamente basta \(F\phi\), nel senso che nel mio modello voglio raggiungere uno stato in cui \(\phi\) è vera, e dunque basta dire che prima o poi \(\phi\) è vera. In maniera più strutturata, possiamo esprimerle come \(\phi_1 U \phi_2\), in cui non solo voglio raggiungere uno stato in cui \(\phi_2\) è vera, ma voglio anche che fino a quel momento il sistema si comporti in un certo modo, ovvero che \(\phi_1\) sia vera.

\paragraph{Safety}
Queste invece sono esprimibili con \(G\phi\), dove \(\phi\) è una formula che esprime la proprietà di sicurezza, ovvero che non accada qualcosa che non voglio che accada. Anche qui, dualmente, posso esprimerla con \(R\), ad esempio \(\phi_1 R \phi_2\) esprime che \(\phi_2\) è sempre vera a meno che non diventa vera \(\phi_1\), ovvero voglio dire non fare mai qualcosa di sbagliato a meno che non sia successo qualcosa che mi dice che non è più un problema, ad esempio un processo è terminato.

La cosa non è così semplice però, perché come abbiamo già visto, è possibile avere delle tracce spurie, ovvero percorsi nel grafo di computazione che sono ammesse solo perché nell'astrazione del modello si è perso qualche dettaglio, ma che in realtà non sono percorsi realizzabili nel sistema concreto.

Se voglio esprimere una proprietà una proprietà che voglio sia vera infinitamente spesso posso usare \(G F \phi\).

In questo modo posso filtrare tracce che non rispettano un vincolo di fairness. Una proprietà di \(G F \phi \to \psi\) è detta di strong fairness.

I problemi di decisione si basano sul modello di computazione, che abbiamo formalizzato tramite una struttura di kirpke \(K\) il cui linguaggio \(L(K)\) è l'insieme di tutte le parole che rappresentano esecuzioni del sistema concreto. Ora, dato un formula \(\phi\) di LTL\@, vogliamo decidere se \(K \models \phi\), ovvero se tutte le parole in \(L(K)\) soddisfano \(\phi\). Quindi se definiamo il linguaggio \(L(\phi) = \set{w \st w,0 \models \phi}\), allora il problema di verifica è decidere se \(L(K) \subseteq L(\phi)\).\footnote{Per brevità se sto iniziando da indice \(0\), posso scrivere \(w,0\models \phi\) come \(w\models \phi\).}

La negazione di questa proprietà è \(L(K) \cap \overline{L(\phi)} \neq \emptyset\), ovvero \(L(K) \cap L(\neg \phi) \neq \emptyset\). Quindi se un comportamento del sistema soddisfa la negata dalla formula, ho trovato un controesempio. Determinare se dato un sistema \(K\) esiste un suo comportamento che soddisfa la negata è il problema di model checking.

\chapter{Lezione 1/04}

Vogliamo introdurre un linguaggio per descrivere sistemi che operano con tempo continuo.

Supponiamo di voler progettare un controllore per un passaggio a livello.

Questo sistema ha una parte digitale, che è quella del controllore, e una parte analogia che sono gli attuatori e i sensori.

Possono essere rappresentate da questi transition system.
\todo{Aggiungi immagini}

Se usiamo il modello classico sarebbe semplicemente composizione asincrona con interleaving, ma questo porterebbe a dei comportamenti ammessi dal modello ma non desiderabili.

\todo{Immagine computation tree}

Quello che vorremmo e che il treno mandi il segnale di essere vicino sufficientemente prima che arrivi al passaggio a livello, e quindi è necessario poter specificare che le azioni impieghino un certo tempo per essere eseguite, che certe azioni possono essere abilitate solo dopo un certo tempo o che certe azioni devono impiegare al più un certo tempo per essere eseguite.

\todo{Immagine con tempistiche}

In questo modo, se vengono garantite queste tempistiche, allora il comportamento non desiderabile non è più ammesso, e dunque il modello è più preciso.

\todo{Immagine con tempistiche e computation tree}

\begin{note}{}
  Nonostante abbiamo due vincoli su una transizione sincronizzante, è importante a livello di modellaizone che siano due e su sottomodelli diversi, piuttosto che accorparli in un unico vincolo.

  Questo perché nella fase di modellazione è importante che ogni sottomodello si occupi di descrivere il sottocomponente specifico e le sue caratteristiche, e dunque è importante che ogni sottomodello si occupi solo dei vincoli del componente che sta modellando.

  Il comportamento complessivo del sistema poi dev'essere corretto grazie a proprietà emergenti dalla composizione dei sottomodelli.
\end{note}

Andiamo a introdurre gli automi temporizzati, come estensione degli automi a stati finiti.

Come per gli automi a pila, prendiamo un automa a stati finiti e aggiungiamo una struttura dati, in questo caso prendiamo un automa a stati finiti e aggiungiamo dei cronometri (o clock), che sono delle variabili i cui valori sono numeri reali non negativi, e che può essere scritto solo in un certo modo, ovvero possono essere resettati a zero o lasciati invariati, ma non possono essere assegnati a valori arbitrari.

Possiamo vedere il clck matematicamente come una funzione \(C: \R^+ \to \R^+\) che è in funzione del tempo reale, tale che \(\frac{d x(t)}{t} = 1\).\footnote{L'importante peril valore di questa derivata è che sia una costante positiva uguale per tutti i clock, ovvero che tutti i clock crescano alla stessa velocità. In tal senso si può prendere tale costante come unità di misura del tempo, e dunque si può assumere che sia \(1\) senza perdita di generalità.} Le transizioni possono avere guardie, le quali possono essere vincoli sui clock. Per semplicità di trattazione, consideriamo solo guardie riguardo ai clock, nonostante in linea di principio le guardie viste in precedenza sono ancora ammesse. Come già detto, l'unica azione che le transazioni possono effettuare sul clock è resettarli a zero.

\begin{note}{}
  Questo modelo è aperto a molte estensioni utili in diverse situazioni, come ad esempio guardi che comprendono l'interazione tra diversi clock:
  \begin{itemize}
    \item vincoli triangolari, ad esempio \(x-y \leq c\). Questi vincoli non cambiano l'espressività dei modelli.
    \item possibilità di vincoli con somma o prodotti come \(x+y \leq c\), questi invece aumentano l'espressività dei modelli e la complessita dei problemi di decisione.
    \item possibilità di resettare un clock a un valore diverso da zero, ad esempio \(x := c\). Anche questo non aumenta l'espressività dei modelli o la complessità dei problemi di decisione.
    \item possibilità di resettare un clock per farlo partire da un valore di un altro clock, ad esempio \(x := y\). Questo invece aumenta l'espressività dei modelli o la complessità dei problemi di decisione.
  \end{itemize}

  In generali i vincoli introdotti per i timed automata sono rilassabili per avere più succintezza, ma non espressività senza esplosione della complessità.
\end{note}

I vincoli sui clock \(\phi(X)\) sono:

\[\phi ::= \top \mid x<c\mid x\leq c\mid x > c \mid x \geq c \mid \phi \land \phi\]

Si noti che non è ammessa la negazione (per quello sono aggiunti tutti i possibili confronti di disuguaglianza) e che non è presente la disgiunzione, perché questa è modellata tramite più transizioni avente come guardie i diversi disgiunti che compongono la disgiunzione. Le costanti \(c\) sono invece di un dominio numerabile, ad esempio \(c \in \N\) o \(c \in \Q\), ma non è possibile avere \(c \in \R\) poiché questo porterebbe a problemi di rappresentazione e computazione.

Un'automa temporizzato è quindi definito come
\[TA = \langle\Sigma, \Loc, \Loc_0,C,\delta,F\rangle\]

Dove \(\Sigma\) è l'alfabeto, \(\Loc\) è l'insieme delle locazioni, \(\Loc_0\) è l'insieme delle locazioni iniziali, \(C\) è l'insieme dei clock, \(\delta \subseteq \Loc \times \Sigma \times \phi(C) \times 2^C \times \Loc\) è l'insieme delle transizioni, e \(F\) è l'insieme delle locazioni finali. Da un punto di vista diagrammatico, una transizione è rappresentata come etichetta con \(t, \phi, a\), dove \(t\) è il trigger, ovvero il simbolo di alfabeto che deve essere letto per poter effettuare la transizione, \(\phi\) è la guardia sui clock che deve essere soddisfatta per poter effettuare la transizione, e \(a\) è l'azione che viene effettuata.

\todo{Aggiungi immagine esempio interruttore}

Questo modello manca ancora di qualcosa. Infatti ci permette di impedire che una transizione avvenga prima di un certo tempo o dopo, e che quindi sia abilitata in un certo intervallo di tempo, ma non ci permette di esprimere che una transizione debba avvenire per forza entro un certo tempo, ovvero che se è abilitata allora deve essere effettuata entro un certo tempo.

Vediamo prima però la semantica, ovvero il transition system associato a un automa temporizzato.
Lo stato è ovviamente una tupla \(\langle l, \nu\rangle\) con \(l \in \Loc\) e \(\nu: C \to \R^+\) che è una funzione di valutazione dei clock.

Dobbiamo quindi definire la relazione di transizione tra stati, che può essere dovuta sia da un simbolo di alfabeto che dal passare del tempo. Quindi \(\Act = \Sigma \cup \R^+\), che porta ad avere due tipi di transizioni, quelle discrete etichettate con simboli di \(\Sigma\) da un \(\langle l,\nu \rangle\) a un \(\langle l',\nu' \rangle\), poiché è possibile azzerare alcuni clock e cambiano locazione, e quelle di passaggio del tempo etichettate con un \(d \in \R^+\) da un \(\langle l,\nu \rangle\) a un \(\langle l,\nu + d\rangle\), in cui non cambia la locazione ma i clock aumentano di \(d\).

In particolare per le transizioni discrete, in \(\delta\) deve esserci una transizione \(l \xrightarrow{t, \phi, a} l'\) tale che \(t\) è il simbolo di alfabeto che viene letto, \(\nu\models\phi\), che dobbiamo definire, e un \(\Eff\) che cambia la valutazione, prendendo in input un insieme di clock da resettare \(\lambda\), e dunque \(\Eff(\lambda, \nu) = \nu[\lambda \to 0]\), dove per \(\nu[\lambda \to 0]\) si intende una nuova funzione \(\nu'\) tale che:
\[\nu'(x) = \begin{cases}
    0      & ,x \in \lambda      \\
    \nu(x) & ,\text{ altrimenti}
  \end{cases}\]

Per \(\nu\models\phi\) è semplice, per le condizioni atomiche è sufficiente sostituire \(x\) con \(\nu(x)\) e verificare se la disuguaglianza è soddisfatta, mentre per le congiunzioni è necessario verificare che siano soddisfatte entrambe le formule.

Per quelle di passaggio del tempo, invece, è necessario definire \(\Eff(d,\nu) = \nu + d\), dove per \(\nu + d\) si intende una nuova funzione \(\nu'\) tale che \(\nu'(x) = \nu(x)+d\), \(\forall x \in C\).

\chapter{Lezione 8/04}

Appunti mancanti per computer scarico, recuperare tramite registrazione audio e appunti fabrizio.

\chapter{Lezione 13/04}

Problema dell'interezione di automi di Büchi.
L'esempio classico è il linguaggio \(L=a^\omega\).

Posso costruire due automi di due stati che accettano \(L\) uno che accetta nel primo stato e l'altro nel secondo.
Entrambi accettano \(L\), perché basta passare infinitamente spesso per uno stato accettante.
Ci aspetteremo quindi che l'interezioni degli automi accetti l'intersezione dei linguaggi, ovvero \(L\).
Ma creando l'automa intersezione tramite prodotto cartesiano, otteniamo un automa che non accetta alcuna parola.

Questo perché gli stati accettanti dell'automa intersezione sono quelli in cui entrambi gli automi originali sono in uno stato accettante, ma tali stati non sono raggiungibili.
Quello che si fa è ammettere che l'accettazione possa avvenire in modo asincrono, ovvero si creano due copie dell'automa intersezione, in cui in una copia si considerano accetanti gli stati in cui il primo automa è in uno stato accettante, e nell'altra copia si considerano accettanti gli stati in cui il secondo automa è in uno stato accettante, e poi si aggiungono delle transizioni che permettono di passare da una copia all'altra, in modo che una volta che il primo automa è passato per uno stato accettante, allora si passa alla copia in cui si considerano accettanti gli stati in cui il secondo automa è in uno stato accettante, e così via, in modo che se una parola è accettata da entrambi gli automi originali, allora è accettata anche dall'automa intersezione.

Formalmente, l'intersezione di automi di Büchi è
\[A_1 \Cap A_2 = \langle\Sigma, Q_1 \times Q_2\times\set{1,2}, \delta', Q_1^0 \times Q_2^0\times\set{1}, F_1 \times Q_2\times\set{1}\cup Q_1\times F_2\times\set{2} \rangle \]

con \(\delta'\) definita come segue:

\[\delta'((q_1,q_2,i),\sigma)=\delta_1(q_1,\sigma)\times\delta_2(q_2,\sigma)\times\begin{cases}
    \set{i}   & ,q_i\notin F_i \\
    \set{3-i} & ,q_i\in F_i
  \end{cases}\]

Quindi gli automi di Büchi sono chiusi per intersezione oltre che per unione. Vedremo successivamente che sono chiusi anche per complementazione, ma non sono chiusi per determinizzazione, ovvero non esiste una procedura per costruire un automa di Büchi deterministico equivalente a un automa di Büchi non deterministico.

Prendiamo ad esempio il linguaggio \(L = \Sigma^* 1^\omega\), con \(\Sigma= \set{0,1}\). In altre parole parole che iniziano in qualsiasi modo ma finiscono con soli \(1\). Possiamo facilmente riconoscere questo linguaggio con il seguente automa di Büchi non deterministico:
\todo{Aggiungi immagine}

Si potrebbe pensare di applicare la costruzione tramite sott'insiemi degli automi a stati finiti per provare a costruire un automa di Büchi deterministico equivalente, ma otterrei il seguente automa:
\todo{Aggiungi immagine}

Chiaramente questo automa, indipendentemente da come definiamo chi siano gli stati accettanti, verranno accettate parole che contengono un numero finito di \(0\), che non sono parole di \(L\).

Si potrebbe pensare quindi che il problema sia nel fatto di aver usato la costruzione tramite sott'insiemi e che esista un'altra procedura per costruire un automa di Büchi deterministico equivalente a un automa di Büchi non deterministico. Questo non è così.

\begin{proof}
  Supponiamo per assurdo che esista un DBA \(D\) tale che \(\mathcal{L}(D)=\mathcal{A} = \Sigma^* 1^\omega\)
  Allora
  \[w_0 = 1^\omega \in \mathcal{L}(D) \implies \exists r_0 \in Q^\omega \text{ run su } w_0 \text{ con } inf(r_0) \cap F \neq \emptyset\]

  Per raggiungere lo stato finale deve esistere un \(q_f \in r_0\) tale che \(q_f \in F\) e \(r_0=r_0' q_f r_0''\).
  Essendo però \(\delta_D\) deterministica, allora la parola sarà del tipo \(w_0=w_0'1w_0''\).

  Andiamo a costruire quindi una parola \(w_1\) tale che dopo l'\(1\) che porta il run ad avere \(q_f\) ha uno \(0\). Ovvero \(w_1 = w_0' 1 0 w_1\).

  Per essere accetata il run corrispondente deve avere.
  \todo{completa da registrazioni}
\end{proof}

AUTOMI GENERALIZZATI DI BÜCHI

L'idea è la stessa di quelli di Büchi classici, ma al posto di avere un unica condizione di accettazione, ne ha più di una.

Struttralmente quindi sarà:

\[A = \langle\Sigma,Q,\delta,Q_0,\langle F_0, \ldots,F_t\rangle\rangle\]
Ovvero al posto di avere un unico insieme di stati accettanti, abbiamo una tupla di insiemi di stati accettanti, e diremo che il linguaggio accettato da questo automa è
\[L(A)= \set{w\in\Sigma^\omega}[inf(r) \cap F_i \neq \emptyset, \forall i \in \set{0,\ldots,t}]\]

Tali automi sono stati introdotti solo per semplicità di trattazione, in quanto non sono più espressivi di quelli di Büchi classici.

In particolare possiamo usare la costruzione di intersezione asincrona che abbiamo visto precedentemente, ma essendo che gli stati sono identici, al posto di avere una copia per ogni automa che stiamo intersecando avremo una copia per ogni insieme di stati finali della tupla. In particolare avremo

\[\hat{A} = \langle\Sigma, \hat{Q},\hat{\delta},\hat{Q_0},\hat{F}\rangle\]
Con \(\hat{Q}=Q\times\set{0,\ldots,t}\), \(\hat{Q_0}=Q_0\times\set{0}\), \(\hat{F}=F_0\times\set{0}\) e la funzione di transizione definita come segue:

\[\hat{\delta}((q,i),\sigma) = \delta(q,\sigma) \times \begin{cases}
    \set{i}                & ,q\notin F_i \\
    \set{(i+1) \mod (t+1)} & ,q\in F_i
  \end{cases}\]

Tale \(\hat{\delta}\) fa in modo che all'arrivo in uno stato accettante di \(F_i\) si passi ad una copia diversa. Di conseguenza, essendo che una parola è accettata se e solo se il run corrispondente visita infinitamente spesso stati accettanti, per tornare infinitamente spesso in stati accettanti di \(F_0\) devo passare per tutte le copie, ovvero per stati accettanti di tutti gli insiemi \(F_i\), e dunque la parola è accettata se e solo se è accettata da tutti gli automi di Büchi classici che hanno come stati accettanti i diversi \(F_i\).

\begin{theorem}{}
  \[\forall \psi \in LTL, \exists A_\psi = \langle 2^{AP}, Q, \delta,Q_0,F\rangle \in NBA \st \mathcal{L}(\psi) = \mathcal{L}(A_\psi)\]
\end{theorem}

Prima di poter dimostrare questa cosa dobbiamo intodurre dei concetti.

\begin{definition}{Chiusura di una formula LTL}
  La chiusura di una formula LTL è una funzione \(cl:LTL\to 2^{LTL}\) tale che:
  \begin{enumerate}
    \item \(\phi \in cl(\phi)\)
    \item \(Op\phi'\in cl(\phi)\implies \phi'\in cl(\phi), \forall Op \in \set{\neg,X}\)
    \item \(\phi_1 Op \phi_2 \in cl(\phi) \implies \phi_1, \phi_2 \in cl(\phi), \forall Op \in \set{\land,\lor}\)
    \item \(\phi_1 U \phi_2 \in cl(\phi) \implies \phi_2 \lor (\phi_1 \land X(\phi_1 U \phi_2)) \in cl(\phi)\)
    \item \(\phi_1 R \phi_2 \in cl(\phi) \implies \phi_2 \land (\phi_1 \lor X(\phi_1 R \phi_2)) \in cl(\phi)\)
  \end{enumerate}
\end{definition}

\begin{example}{}
  Sia \(\phi=p\land X(rUq)\).
  La chiusura avrà i seguenti elementi:
  \begin{itemize}
    \item \(p\land X(rUq)\) per la regola 1
    \item \(p\) e \(X(rUq)\) per la regola 3
    \item \(rUq\) per la regola 2
    \item \(q \lor (r \land X(rUq))\) per la regola 4
    \item \(q\) per la regola 3
    \item \(r \land X(rUq)\) per la regola 3
    \item \(r\) per la regola 3
    \item \(X(rUq)\) per la regola 2
  \end{itemize}

  Questo esempio ci fa notare che le regole 4 e 5 portano ad aggiungere anche formule similmente alla regola 2.

  Si può facilmente vedere che \(2\abs{\phi}\geq \abs{cl(\phi)}\).
\end{example}

La chiusura è una nozione puramente sintattica, infatti tratta ugualmente cose come \(\land\) e \(\lor\). Di fatto la chiusura non è altro che l'insieme delle sottoforumle con l'aggiunta del one-step unfolding per \(U\) e \(R\).

\begin{definition}{Atomi di formule LTL}
  Data una formual \(\phi\) definiamo gli atomi di \(\phi\) come l'insime \(ATM(\phi)\subseteq 2^{cl(\phi)}\) tale che, un certo \(A \in ATM(\phi)\) se e solo se:
  \begin{itemize}
    \item \(A\subseteq cl(nnf(\phi))\)
    \item \(\forall p \in AP, p\in A \oplus  \neg p \in A\), ovvero \(p\in A\iff \neg p \notin A\)
    \item \(\phi_1 \land \phi_2 \in A \implies \phi_1, \phi_2 \in A\)
    \item \(\phi_1 \lor \phi_2 \in A \implies \phi_1 \in A \text{ o } \phi_2 \in A\)
    \item \(\phi_1 R \phi_2 \in A \implies \phi_2 \land (\phi_1 \lor X(\phi_1 R \phi_2)) \in A\)
    \item \(\phi_1 U \phi_2 \in A \implies \phi_2 \lor (\phi_1 \land X(\phi_1 U \phi_2)) \in A\)
  \end{itemize}
\end{definition}

Notare che l'operatore \(X\) non è presente nelle condizioni dell'atomo, questo perché la semantica di \(X\) necessità di parlare di più punti della parola, e dunque non è possibile definire una condizione che riguarda solo un punto della parola, come invece è possibile fare per gli altri operatori.

\begin{definition}{Sequenza di Hintikka}
  Una sequenza di Hintikka per una formula \(\phi\) è una parola infinita di atomi di \(\phi\) ovvero \(\alpha \in (ATM(\phi))^\omega\) tale che:
  \begin{enumerate}
    \item \(\phi \in \alpha_0\)
    \item \(\forall i, X\phi' \in \alpha_i \implies \phi' \in \alpha_{i+1}\)
    \item \(\forall i, \phi_1U\phi_2 \in \alpha_i \implies \exists j \in \N\st\phi_2\in\alpha_{i+j}\)
  \end{enumerate}
\end{definition}

\begin{theorem}{}
  \[\forall \phi \in LTL, Sat(\phi) \iff \exists \alpha \in (ATM(\phi))^\omega\st\alpha \text{ è una sequenza di Hintikka per } \phi\]
\end{theorem}
\begin{proof}
  VEDI REGISTRAZIONE
\end{proof}

Possiamo finalmente dimostrare il teorema di traduzione da LTL a NBA\@.
Per fare ciò anremo a costruire un automa di Büchi \(A_\phi\) che accetta esattamente le sequenze di Hintikka per \(\phi\), e dunque, essendo che una parola soddisfa \(\phi\) se e solo se esiste una sequenza di Hintikka per \(\phi\) che è uguale alla parola, allora \(A_\phi\) accetterà esattamente le parole che soddisfano \(\phi\).

Sia quindi \[A_\phi = \langle 2^{AP}, Q, \delta, Q_0, \hat{F} = \langle F_0, \ldots, F_t\rangle\rangle\]
Dove \(Q = ATM(\phi)\), \(Q_0 = \set{A \in ATM(\phi) }[\phi \in A]\), la funzione di transizione \(\delta\) si assicurerà che varrà la seconda condizione delle sequenze di Hintikka, ovvero
\[\delta(A,\sigma)=\begin{cases}
    \set{A'\in ATM(\phi)}[\forall X\psi \in A, \psi \in A'] & , \sigma = A \cap AP \\
    \emptyset                                               & , \text{ altrimenti}
  \end{cases}\]
mentre gli stati di accettazione saranno quelli che soddisfano la terza condizione delle sequenze di Hintikka, creando un insieme di stati di accettazione per ogni formula del tipo \(\phi_1 U \phi_2\) presente nella chiusura, ovvero
\[\hat{F} = \set{\set{A\in ATM(\phi)}[\phi_1 U \phi_2 \in A \implies \phi_2\in A]}[\phi_1 U \phi_2 \in cl(\phi)]\]

Tale costruzione ci permette di accettare solo parole che passano infinitamente spesso per stati in cui ogni until è soddisfatto localmente, ovvero \(\phi_1 U \phi_2\) è vera perché \(\phi_2\) è vera.


\begin{example}{}
  Sia \(\phi= F p \land q U r\).
  \[cl(\phi) = \set{\phi, Fp,qUr, r\lor q \land X(q\cup r),p\lor XFp, r,p,q,X(qUr),xFp}\]

  Fortunatamente non tutte le combinazioni di questi elementi sono atomi validi.
  Possiamo partire dagli stati iniziali, ovvero dagli atomi che contengono \(\phi\),

  INSERIRE IMMAGINI ATOMI
\end{example}

\chapter{Lezione 15/04}

In termini di teoria dei linguaggi questi automi ci permettono di esprimere linguaggi temporizzati, in cui gli elementi delle parole sono dette segnali, ovvero coppie di un simbolo di alfabeto e un istante di tempo in cui tale simbolo viene emesso. In tal senso i run di un automa temporizzato su una traccia temporizzata è una sequenza di transizione del transition system come in figura

\todo{Vedi slide, aggiungi immagini}

Ovviamente nei percorsi temporizzati per ogni percorso tra due transizioni discrete ci sono infiniti altri sottopercorsi che sezionano la transizione temporale in sotto transizioni con ritardi la cui somma è la stessa della transizione originale.

In particolare per ogni transizione con ritardo \(d\), esiste una coppia di transizioni con ritardo \(d_1\) e \(d_2\) equivalente per ogni coppia tale che \(d=d_1+d_2\).

Per questi automi è possibile avere transizioni discrete solo con stimoli esterni, per fare in modo di avere una transizione spontanea dobbiamo aggiungere le \(\tau-\)transizioni, ovvero aggiungere un simbolo di alfabeto \(\tau\) che indica una transizione che avviene senza stimolo esterno. A differenza delle \(\epsilon-\)transizioni negli automi a stati finiti, le \(\tau-\)transizioni aumentano il potere espressivo degli automi temporizzati.

Questo lo possiamo vedere nell'automa temporizzato con \(\tau-\)transizioni sottostante.

Nessun automa temporizzato senza \(\tau-\)transizioni può riconoscere lo stesso linguaggio di questo automa, perché la transazione spontanea azzera il clock, permettendo di accettare tutte le parole del tipo \((a,t_1)(a,t_2)\ldots(a,t_n)\ldots\) con \(t_i \in \N\). Questo perché a ogni istante di tempo intero il clock può essere azzerato, potendoli quindi scandire tutti. Non avendo la possibilità di azzerare il clock in maniera spontanea invece non sarebbe possibile scandire il clock, perché potrei solamente valutare che c'è un intero specifico di distanza fra due clock, ma non di azzerarli in modo da poter scandire tutti gli interi.

Dobbiamo ancora capire come forzare l'avvenimento di una transizione. Queste \(\tau-\)transizioni per ora ci permettono solo di dare la possibilità che avvenga, non forzano l'automa a effettuarla.

Per fare ciò, si introduce una funzione detta di invariante \(Inv: \Loc \to \Phi(C)\) che associa un vincolo sui clock a ogni locazione, la cui semantica è di limitare il tempo che è possibile trascorrere in una locazione, ovvero se \(Inv(l) = x \leq c\) allora è possibile rimanere in \(l\) solo per un tempo al più \(c\). Il senso è forzare che ogni percorso che non rispetti l'invariante termini, ovvero che ogni percorso che non la rispetta sia considerata un'esecuzione spuria del sistema.

Ovviamente non è tutto rose e fiori, perché l'introduzione di questa invariante può portare all'interruzione del sistema, che in linea di principio può introdurre condizioni simili al deadlock, detta timelock, in cui il sistema deve evolvere perché sta per andare a incontro a una deadline dovuta a una invariante ma tutte le transizioni sono disabilitate e quindi non può fare niente.

Dobbiamo aggiornare la semantica per introdurre questa invariante, in particolare inserendo transizioni di passaggio del tempo solo se tale passaggio soddisfa l'invariante, ovvero se \(\nu + d \models Inv(l)\). Per quando riguarda le transizioni discrete invece, è invece necessario assicurarsi che la valutazione del clock dopo l'effetto della transizione soddisfi l'invariante, ovvero che \(\nu' \models Inv(l')\).

L'aggiunta delle invarianti ci permette quinidi di modellare le durate delle azioni, oltre che hai ritardi già espressi dalle guardie.

Per quanto riguarda la composizione di automi temporizzati, avendo considerato il caso senza variabili, possiamo utilizzare il modello di sincronizzazione con handshaking, in cui le transizioni sincronizzanti sono quelle che hanno lo stesso simbolo di alfabeto, mentre le altre transizioni sono asincrone.

Andiamo a vedere nel dettaglio qualche problema di modellazione che può avvenire negli automi temporizzati.

Un primo problema è la time convergence, che è inevitabile, ma che può essere risolto in fase di verifica.
Un secondo problema è quello del timelock, che è introdotto dalle invarianti, e ne abbiamo già parlato, ma approfondiremo.
L'ultimo è quello di zenoness, in cui infinite azioni avvengono un intervallo di tempo finito, e questo può avvenire essendo che le transizioni discrete sono istantanee e in particolare quelle spontanee non necessitano di uno stimolo esterno. Questo è un problema perché ametterle porterebbe a un sistema con una velocità di esecuzione arbitrariamente alta.

Andiamoli a definire nel dettaglio.

Si ha time convergence quando esiste un percorso infinito il cui tempo converge verso un valore finito.

Questo è facile vedere che possa avvenire, ad esempio con l'automa in figura.

Chiamando il percorso \(\pi\) possiamo definire il tempo di esecuzione di \(\pi\) come
\[ExecTime(\pi) = \sum_{i \geq 1} ExecTime(d_i) = \sum_{i \geq 1} \frac{1}{2^i} = 1\]

Se \(ExecTime(\pi)<\infty\) allora \(\pi\) è un percorso convergente. Tali percorsi possono essere ignorati se il numero di transizioni discrete in \(\pi\) è finito, perché in tal caso è possibile accorpare le infinite transizioni temporali in un'unica transizione temporale con ritardo pari al limite della somma dei ritardi, e dunque non si perderebbe alcuna parola accettata dall'automa.

Dato uno stato di partenza \(s\) ci interesseremo solo dell'insieme dei percorsi che partono da \(s\) e che sono divergenti, ovvero:
\[Path_{div}(S) = \set{\pi}[\pi_0 = s \land ExecTime(\pi) =\infty]\]

Definiamo invece un timelock come uno stato \(s\) tale che \(Path_{div}(s) = \emptyset\), ovvero non esiste alcun percorso divergente che parte da \(s\). Questo non vuol dire che da \(s\) non parta alcun percorso, o che il sistema non evolva, ma che tutti i percorsi che partono da \(s\) sono convergenti. Ad esempio, nell'automa in figura

L'automa può evolvere se non in percorsi convergenti.

Per quanto riguarda la zenoness, la situazione è più complessa. Un percorso \(\pi\) è zeno se \(ExecTime(\pi)<\infty\) e il numero di transizioni discrete in \(\pi\) è infinito. Un automa temporizzato è non-zeno se non ammette percorsi zeno, ovvero se ogni percorso con un numero infinito di transizioni discrete è divergente.

Purtroppo la verifica di presenza di zenoness è un problema difficile, ma conosciamo una condizione sufficiente per la non-zenoness.

In particolare, se in ogni ciclo semplice dell'automa temporizzato esiste un clock \(x\) che rispetta le seguenti proprietà:
\begin{itemize}
  \item \(x\) è resettato in almeno una transizione del ciclo
  \item in ogni valutazione del ciclo, esiste un intero positivo \(c\) tale che \(x \geq c\) è una guardia di almeno una transizione del ciclo
\end{itemize}
allora l'automa temporizzato è non-zeno. Infatti ognuno di questi cicli richiede un tempo di esecuzione di almeno \(c\) per essere attraversato. Il percorso indotto da tale ciclo nell'LTS associato non è necessariamente un ciclo, perché gli stati possono essere diversi a causa della valutazione dei clock, ma considerando solo le locazioni, necessariamente dovrà contentere dei cicli. In particolare tali cicli devono essere attraversati un numero infinito di volte, essendo che il numero di cicli semplici è finito, e dunque il tempo di esecuzione del percorso è infinito, poiché sarà almeno \(c\) per ogni ciclo attraversato, e dunque il percorso non è zeno.

\chapter{Lezione 20/04}

A questo punto dobbiamo capire come usare questi automi per risolvere problemi di model checking e soddisfacibilità.

In particolare, data una struttura di kripke dobbiamo capire se questa soddisfa una formula \(\psi\). Questo equivale a dire che
\(\mathcal{L}(k)\subseteq \mathcal{L}(\psi) = \mathcal{A_\psi}\). Questo però equivale a richiedere che tutti i percorsi di \(k\) siano accettati da \(\mathcal{A_\psi}\).

Per semplicità, piuttosto di considerare che tutti i percorsi di \(k\) siano accettati da \(\mathcal{A_\psi}\), possiamo considerare il modello della negata e verificare che l'intersezione sia vuota.

Ricordiamo che la struttura di kripke \(k\) è definita come \(k = \langle W, R, AP, \lambda, w_0\rangle\) con \(W\) insieme dei mondi possibili o stati, \(R \subseteq W\times W\) relazione di transizione, \(AP\) insieme di proposizioni atomiche, \(\lambda: W \to 2^{AP}\) funzione di etichettatura che associa a ogni stato un insieme di proposizioni atomiche vere in quello stato, e \(w_0\) mondo iniziale.

A quel punto cotruire un automa risulta molto facile considerando \(A_k = \langle \Sigma=2^{AP}, Q = W, \delta, q_0 = w_0,  F=W\rangle\) con
\[\delta(w,\sigma) = \begin{cases}
    \set{w' \in W}[R(w,w')] & , \text{ se } \sigma = \lambda(w) \\
    \emptyset               & , \text{ altrimenti}
  \end{cases}\]

Tale automa accetta tutte e sole le parole del linguaggio della struttura di kripke \(k\), e dunque se l'intersezione di prima risulta essere \(\mathcal{L}(A_k) \cap \mathcal{L}(A_{\neg\psi})\), ovvero \(A_{k,\neg\psi} = A_{\neg\psi} \cap \mathcal{L}(A_k)\).
Quindi \(k\models \psi\) se e solo se \(A_{k,\neg\psi} = \emptyset\).

Questo va bene per il model checking, per la soddisfacibilità invece è sufficiente costruire \(A_\psi\) e verificare che \(\mathcal{L}(A_\psi)\neq \emptyset\). Per la validitò uguale ma usando \(A_{\neg\psi}\) e controllando che \(\mathcal{L}(A_{\neg\psi}) = \emptyset\).

Quindi l'ultima cosa che ci serve è un algoritmo per verificare se il linguaggio di un automa di Büchi è vuoto o meno.

Consideriamo per ora un NFA dato da \(A = \langle \Sigma,Q,\delta,Q_0,F\rangle\) e cerchiamo di capire se \(\mathcal{L}(A) = \emptyset\).

Consideriamo inizialmente l'esempio in figura.

\todo{Aggiungi immagine}

Intuitivamente per capire se è vuoto o meno dobbiamo vedere se esiste un percorso da uno stato iniziale a uno stato di accettazione, ovvero che uno stato finale sia raggiungibile dallo stato iniziale.

Trovato un percorso infatti posso sembre vedere quali archi abbiamo attraversato e usarne le etichette per sintetizzare la parola accettata.

Nel'esempio di prima \(S_0 S_1 S_2\) è un percorso accettante, le cui etichette sono \(bb\). Questo ci dice che \(bb\in \mathcal{L}(A)\).

Un modo banale è di usare un normale algoritmo di reachability sul grafo trasposto a partire dagli satati finali, e vedere se riesco a raggiungere uno stato iniziale, verificabile con un iterazione sugli stati raggiunti.

Esiste però un algoritmo non in NLOGSPACE, ovvero non deterministico a cui basta spazio logaritmico.

TODO PULISCI ALGORITMO E USA AMBIENTE APPROPRIATO

\begin{enumerate}
  \item Guess \(s \in Q_0\)
  \item Guess \(\sigma \in \Sigma\)
  \item Guess \(s' \in \delta(s,\sigma)\)
  \item If \(s' \in F\) then
  \item return yes
  \item else
  \item s \(\gets t\)
  \item goto 2
\end{enumerate}

Questo è non deterministico perché ci sono dei guess, valuiamo il perché è LOGSPACE\@.

Qui abbiamo 3 variabili \(s, \sigma, s'\) ma ognuna di queste può rappresentare un elemento di \(Q\) o \(\Sigma\) in spazio logaritmicol, in particolare \(2\log(\abs{Q}) + \log(\abs{\Sigma})\).

Sapendo che lo spazio è limitato, possiamo aggiungere la condizione di terminazione per la procedura contando il numero di archi attraversati, se supera il numero di stati sono in un ciclo sicuramente.

\begin{enumerate}
  \item count \(\gets 0\)
  \item Guess \(s \in Q_0\)
  \item Guess \(\sigma \in \Sigma\)
  \item Guess \(s' \in \delta(s,\sigma)\)
  \item If \(s' \in F\) then
  \item   return yes
  \item else if count < \(\abs{Q}\)
  \item   s \(\gets t\)
  \item goto 2
  \item else
  \item   return no
\end{enumerate}

Questa nuova variabile è ancora logaritmica.

Dobbiamo ora estendere agli automi di Büchi, in cui non voglio semplicemente raggiungere uno stato, ma voglio raggiungere uno stato infinite volte.

Possiamo estendere l'algoritmo precedente facendo guess dello stato da raggiungere, per potermelo ricordare, e poi verificando che ci ritorno.

\begin{enumerate}
  \item Guess \(s \in Q_0\)
  \item Guess \(\sigma \in \Sigma\), Guess \(s' \in \delta(s,\sigma)\)
  \item \(s'\gets t\)
  \item if \(s' \in F\) then
  \item   Guess \(\sigma \in \Sigma\), Guess \(s \in \delta(s',\sigma)\)
  \item   if \(s = s'\) then
  \item     return yes
  \item   else
  \item     goto 5
  \item else
  \item  goto 2
\end{enumerate}

Questo algoritmo è ancora in NLOGSPACE\@, e prima trova un percorso verso uno stato finale, e poi verifica che esiste un ciclo che torna a quello stato finale, dunque almeno un modo per passarci infinitamente spesso esiste.

Per farlo in maniera direttamente deterministica calcolando le componenti fortemente connesse del grafo, e verificando se esiste uno stato finale che è in una componente fortemente connessa. Di quelli che hanno questa proprietà, bisogna verificare che siano raggiungibili da uno stato iniziale, e se è così allora il linguaggio è non vuoto, altrimenti è vuoto.

Per gli automi di CO-Büchi invece, è sufficiente escludere i nodi di rifiuto nel calcolo delle componenti fortemente connesse, in modo da individuare cicli che non passano mai per stati di rifiuto, e verificare che tali cicli siano raggiungibili dallo stato iniziale, eventualmente passando anche per stati di rifiuto (poiché solo un numero finito di volte).

Andiamo quindi a dimostrare che con gli automi di Büchi accettano tutti e soli i linguaggi omega-regolari.

Partiamo dal teorema sugli automi a stati finiti.

\begin{theorem}
  Un linguaggio \(L\) è regolare se e solo se \(\exists NFA A \st \mathcal{L}(A) = L\)
\end{theorem}

Questo è facile da vedere ed è dato per scontato dal corso di EDIT.

\begin{theorem}
  Un linguaggio \(L\) è \(\omega-\)regolare se e solo se \(\exists NBA A \st \mathcal{L}(A) = L\)
\end{theorem}

\begin{proof}
  \begin{description}
    \item[(\(\implies\))] \(L\) \(\omega-\)regolare \(\implies \exists \gamma = \alpha_1\beta_1^\omega+\alpha_2\beta_2^\omega+\cdots + \alpha_n\beta_n^\omega \st \mathcal{L}(\gamma) = L\). Ovviamente ogni \(\alpha_i,\beta_i\) sono regolari, quindi dal teorema precedente esistono \(DFA A_{i}, B_{i}\) tali che \(\mathcal{L}(A_i) = \alpha_i\) e \(\mathcal{L}(B_i) = \beta_i\).
      Cotruiamo quindi per ogni componente \(i\) un automa di Büchi \(C_i= \langle \Sigma, Q_A \cup Q_B, \delta, \set{q_{0_A}}, \set{q_{0_B}}\rangle\) con \(\delta\) definita come:
      \[\delta(q,\sigma)=\begin{cases}
          \set{\delta_A(q,\sigma)}          & ,q \in Q_A \land \delta_A(q,\sigma)\notin F_A \\
          \set{\delta_A(q,\sigma), q_{0_B}} & ,q \in Q_A \land \delta_A(q,\sigma)\in F_A    \\
          \set{\delta_B(q,\sigma)}          & ,q \in Q_B \land \delta_B(q,\sigma)\notin F_B \\
          \set{\delta_B(q,\sigma), q_{0_B}} & ,q \in Q_B \land \delta_B(q,\sigma)\in F_B    \\
        \end{cases}\]

      Tale costruzione fa si che ogni componente si comporti come la parte regolare \(\alpha_i\), fino a non-deterministicamente decidere di passare alla parte regolare \(\beta_i\), che è estesa per essere un ciclo, in modo da passare infinitamente spesso per lo stato \(q_{0_B}\) che è l'unico stato accettante, e dunque accettare esattamente le parole di \(\alpha_i\beta_i^\omega\). A questo punto mi è sufficiente fare l'unione di tutti questi automi \(C_i\) per ottenere un automa di Büchi che accetta esattamente le parole di \(\gamma\), ovvero di \(L\).
    \item[(\(\impliedby\))] Consideriamo il NBA \(N = \langle \Sigma, Q, \delta, Q_0, F\rangle\) con \(F = \set{q_0\ldots,q_k}\) stati finali di \(N\). Per \(i \in \set{1,\ldots,k}\) definiamo \(A_i = \langle \Sigma, Q, \delta, Q_0, \set{q_i}\rangle\) e \(B_i = \langle \Sigma, Q, \delta, \set{q_i}, \set{q_i}\rangle\). Interpretando ognuno di questi automi come NFA otteniamo che i linguaggi accettati da \(A_i\) e \(B_i\) sono regolari per il teorema precedente. Inoltre, possiamo separare il nostro \(N\) in \(k\) sottoautomi, ognuno dei quali con \(F_i = \set{q_i}\). Ovviamente il linguaggo dell'unione di questi sottoautomi è uguale a quello di \(N\). Ognuno di questi sottoautomi però ha linguaggio \(L(N_i) = L(A_i)\cdot L(B_i)\omega\). Dunque \(L(N) = \bigcup_{i=1}^k L(A_i)\cdot {L(B_i)}^\omega\), ovvero è \(\omega-\)regolare.
  \end{description}
\end{proof}

Quindi i NBA sono sufficientemente espressivi per accettare tutti e soli i linguaggi \(\omega-\)regolari, a differenza di come abbiamo visto per LTL. Lo abbiamo osservato con l'esempio dell'espressione regolre  che accetta le parole che hanno P vera in posizione pari. Questo è facilisimo però da fare con un automa di Büchi, banalmente

P
->
SF  N
<-
*

Vediamo ora un elenco di diverse condizioni di accettazioni per automi che caratterizzano diversi tipi di linguaggi e proprietà di essi. Considerato \(\rho\) un run di un automa avremo:

\begin{itemize}
  \item Accettazione al finito, \(last(\rho) \in F\)
  \item Accettazione di Buchi, \(\inf(\rho) \cap F \neq \emptyset\) o CO-Buchi, \(\inf(\rho) \subseteq F = \emptyset\)
  \item Generalized Buchi, \(\bigwedge_i \inf(\rho) \cap F_i \neq \emptyset\) o Generalized CO-Buchi, \(\bigvee_i \inf(\rho) \cap F_i = \emptyset\)
  \item Accettazione di Muller, abbiamo \(F\subseteq 2^Q\), \(\inf(\rho) \in F\). Questa è più simile a quella a finito che alle altre, perché per accettare richiediamo una condizione specifica per il comportamento all'infinito, mentre in quelli di buchi è sufficiente che uno stato sia attraversato infinitamente spesso il resto non mi interessa. Questa è molto più costosa però, perché non mi basta usare raggiungibilità, ma devo verificare com'è fatto il run, e verificare se la sua parte infinita è in \(F\).
  \item Accettazione di Rabin, abbiamo \(F = {(2^Q\times 2^Q)}^*\), quindi \(F = \langle (G_i,B_i),\ldots,(G_k, B_k)\rangle\), \(\bigvee_i \inf(\rho) \cap G_i \neq \emptyset \land \inf(\rho) \cap B_i = \emptyset\)
  \item Accettazione di Street, abbiamo \(F = {(2^Q\times 2^Q)}^*\), quindi \(F = \langle (G_i,B_i),\ldots,(G_k, B_k)\rangle\), \(\bigwedge_i \inf(\rho) \cap G_i \neq \emptyset \lor \inf(\rho) \cap B_i = \emptyset\)
\end{itemize}

Le ultime 3 nel caso non deterministico si possono ridurre a quelle di Büchi, con un costo non indifferente, mentre per il determinismo non è possibile trovare una riduzione a quelle di Büchi, e dunque sono più espressive.

\chapter{Lezione 22/04}

Tool SPIN che usa Program Graph e LTL. Spin in realtà è più vicino a un transpiler, perchè dato un file con il modello non fa la verifica ma crea un file C che incorpora il modello e la proprietà da verificare con il verificatore, che va compilato e lanciato per avere il risultato della verifica. Spin permette anche la simulazione del modello, sia in maniera guidata dall'utente, che in maniera casuale, che guidata da una traccia di errore.

Per specificare il modello in SPIN si usa un linguaggio di specifica detto Promela. Un modello in Promela è una collezione di \q{proctype} che sono prototipi di processo che vengono istanziati ed eseguiti in parallelo.

\todo{inserire immagine di esempio}

Ci sono diverse istruzioni nel linguaggio che sono considerate atomiche, ovvero che verranno eseguite senza interruzioni, e dunque mappate su una singola transizione nel transition system. I conditional statement e i repetition statement invece sono particolari poiché vogliamo modellare anche non determinismo. In particolare i costrutti "if" e "do" sono un elenco di guardie e azioni, e se più di una guardia è vera allora viene scelto non deterministicamente quale azione eseguire. L'unica differenza tra i due è che "if" viene eseguito una sola volta, mentre "do" viene eseguito finché non raggiunge un break o un goto.

\todo{inserire immagine su if e do}

Una differenza interessante da notare rispetto al C è che nel linguaggio C ogni istruzione è un espressione, ovvero ha un valore che posso assegnare a una variabile. In Promela invece ogni espressione è un istruzione, nel senso che un confronto booleano è un istruzione che è eseguibile se è vera e non eseguibile se è falsa. In particolare, la verifica per le condizioni è fatta sull'eseguibilità dell'istruzione, e quindi è possibile usare ogni istruzione come guardia, e non solo confronti booleani.

Per le formule LTL invece, SPIN usa la sintassi modale per rappresentare gli operatori temporali, ovvero \(\square\) per \(G\), \(\diamond\) per \(F\), \(U\) per \(U\) e \(V\) per \(R\), con \(\square\) rappresentato come [] e \(\diamond\) rappresentato come <>. \todo{aggiungi listing in line}.

Tipi standard più Channel e PID per i canali e i process ID. Approfondisci con slide.

  [ESEMPI IN SPIN]

\chapter{Lezione 27/04}

Passiamo ora dalle logiche temporali lineari a quelle ramificate, in particolare con \(CTL^*\) (full computation tree logic) che contiene LTL ma permette di esprimere anche proprietà che non sono esprimibili in LTL\@.

In particolare vedremo il frammento CTL, che è un sottoinsieme di \(CTL^*\) che però è più facile da verificare, di cui vorremo risolvere il model checking. In particolare vedremo inizialmente un algoritmo diretto per il model checking di CTL, e poi vedremo un alogoritmo più efficiente che usa una struttura dati specifica chiamata OBDD\@.

In particolare questa logica aggiunge dei quantificatori sui percorsi, \(A\) e \(E\), che indicano rispettivamente che una proprietà è vera su tutti i percorsi o su almeno un percorso.

\begin{definition}{Grammatica \(CTL^*\)}
  La grammatica di \(CTL^*\) è formata da due grammatiche dette grammatica delle formule di stato \(\phi\) e grammatica delle formule di percorso\(\psi\), definite come segue:
  \[\phi ::= p \in AP \mid \neg \phi \mid \phi \land \phi \mid \phi \lor \phi \mid A\psi \mid E\psi\]
  \[\psi ::= \phi \mid \neg \psi \mid \psi \land \psi \mid \psi \lor \psi \mid X\psi \mid \psi U\psi \mid \psi R\psi\]

\end{definition}
Possiamo qundi vedere \(\psi\) come \(LTL\) su \(\phi\). Per la soddisfacibilità in \(LTL\) abbiamo detto che per \(\psi \in LTL\) \(k \models_{LTL} \psi\) se ogni percorso della struttura di kripke soddisfa la proprietà. Possiamo modellarlo anche in \(CTL^*\) con \(k \models_{CTL^*} A\psi\), ovvero \(\psi\) è vera su tutti i percorsi.

Passando alla semantica per \(\psi\) è praticamente a LTL se non per il caso atomico in cui va generalizzato da \(p\) a \(\phi\), mentre per la semantica per \(\phi\) è più complessa.

Per i casi base e per i connettivi booleani è abbastanza intuitiva, mentre per i quantificatori sui percorsi si ritorna a \(\psi\).

\begin{definition}{Semantica \(CTL^*\)}
  Data una struttura \(k = \langle AP, W, R, w_0, r \rangle\) con \(AP\) insieme di proposizioni atomiche, \(W\) insieme dei mondi possibili o stati, \(R \subseteq W\times W\) relazione di transizione, \(w_0\) mondo iniziale e \(r: W \to 2^{AP}\) funzione di etichettatura che associa a ogni stato un insieme di proposizioni atomiche vere in quello stato, e dato uno stato \(w \in W\) di partenza (\(w_0\) se omesso), diremo che:
  \begin{itemize}
    \item \(k,w \models p\) se e solo se \(p \in r(w)\)
    \item \(k,w \models \neg \phi\) se e solo se \(k,w \not\models \phi\)
    \item \(k,w \models \phi_1 \land \phi_2\) se e solo se \(k,w \models \phi_1\) e \(k,w \models \phi_2\)
    \item \(k,w \models \phi_1 \lor \phi_2\) se e solo se \(k,w \models \phi_1\) o \(k,w \models \phi_2\)
    \item \(k,w \models A\psi\) se e solo se \(\forall \pi \in Paths(k,w), k,\pi \models \psi\)
    \item \(k,w \models E\psi\) se e solo se \(\exists \pi \in Paths(k,w), k,\pi \models \psi\)
  \end{itemize}
  Importante notare che fino ai quantificatori sui percorsi la semantica è definita su stati, quindi non ha alcuna nozionne temporale, mentre nei quantificatori si passa alla semantica LTL\@, che è definita sui percorsi, e dunque ha una nozione temporale. Normalmente a LTL basterebbe \(\pi\), ma essendo che in una formula di percorso posso tornare a una formula di stato è necessario mantenere la struttura \(k\) per poter valutare le formule di stato che possono essere presenti in \(\psi\).

  In particolare, per la semantica di \(\psi\) abbiamo esattamente la stessa semantica di LTL\@, con l'unica differenza che
  \[k, \pi, i\models \text{ se e solo se} k,(\pi)_i\models \phi\]
\end{definition}

Vedimo ora il frammento CTL di \(CTL^*\), che è un sottoinsieme di \(CTL^*\) in cui ogni operatore temporale è preceduto da un quantificatore sui percorsi.

Quello che cambia è la grammatica di \(\psi\), che non è l'intera grammatica di LTL\@, ma solo:
\[\psi ::= X\phi \mid \phi U\phi \mid \phi R\phi\]

In questo modo non c'è nesting di operatori temporali, e dunque ogni operatore temporale è preceduto da un quantificatore sui percorsi, e dunque ogni formula di percorso è del tipo \(X\phi\), \(\phi U \phi\) o \(\phi R \phi\).

Una formulazione alternativa e la prima proposta usa un'unica grammatica:
\[\phi ::= p \in AP \mid \neg \phi \mid \phi \land \phi \mid \phi \lor \phi \mid AX\phi \mid EX\phi \mid \phi AU \phi \mid \phi AR \phi \mid \phi EU \phi \mid \phi ER \phi\]

Questa sintassi esplicità è più prolissa ma è usata in molti software perché il problema di model checking di CTL è risolto risolvendo dei sottoproblemi per le sottoformule di questa forma e poi combinando.

Essendo molto prolissa, se ne utilizza una versione in cui si rimuovono quelli non necessari, ad esempio \(\land\) in presenza di \(\lor\) e \(\neg\).
Chiaramente questo vale per ogni coppia di operatore duale, ad esempio \(A\) e \(E\), \(U\) e \(R\), \(X\) e \(X\).

Dunque \(\neg AX\phi \equiv E\neg X \phi \equiv EX \neg \phi\), quindi l'operatore \(AX\) è ridondante.
Similmente \(\neg A \phi_1 U \phi_2 \equiv E \neg (\phi_1 U \phi_2) \equiv E \neg \phi_1 R \neg \phi_2\), e di conseguenza \(AU\) è ridondante. Ugualemente \(AR\) è ridondante. Quindi a livello algoritmico è sufficiente considerare solo il \q{frammento esistenzaile} di CTL\@, ovvero quello che contiene solo \(E\) come quantificatore.

Consideriamo inoltre, \(\phi_1 R \phi_2\). Supponiamo che \(\phi_1\) non sia mai vera, allora \(\phi_1 R \phi_2\) è vera se e solo se \(\phi_2\) è sempre vera, ovvero \(G\phi_2\). Altrimenti \(\phi_1\) dev'essere sempre vera fino a quando \(\phi_2\) non diventa vera, e dunque \(\phi_1 R \phi_2\) è vera se e solo se \(\phi_1 U \phi_2\). Dunque \(\phi_1 R \phi_2 \equiv G \phi_2 \lor \phi_2 U (\phi_1 \land \phi_2)\).

Formalmente \(\phi_1 R \phi_2\) è definita come:
\[\forall j \geq i, (\pi, j \models \phi_2 \lor \exists k \geq i, \pi, k \models \phi_1 \land \phi_2)\]

Il primo disgiunto è esattamente \(G\phi_2\), mentre il secondo disgiunto è esattamente \(\phi_2 U (\phi_1 \land \phi_2)\).
Questo vale a livello di LTL, ma a livello di CTL dobbiamo vederlo usando solo \(\lor\) e \(E\).
\[E \phi_1 R \phi_2 \equiv E(G\phi_2 \lor \phi_2 U(\phi_1 \land \phi_2))\]

Essendo che \(E\) distribuisce su \(\lor\) abbiamo proprio
\[E \phi_1 R \phi_2 \equiv E G\phi_2 \lor E \phi_2 U(\phi_1 \land \phi_2)\]

Dunque la logica che andremo veramente a risolvere è la seguente:
\[\phi::= p \in AP \mid \neg \phi \mid \phi \lor \phi \mid EX \phi \mid E\phi U \phi \mid EG \phi\]

Quello che andremo a fare è un algoritmo bottom-up, in cui risolviamo le sottoformule e andremo ad etichettare il nostro modello con delle etichette che indicano quali formule sono soddisfatte in ogni stato.

Prendiamo ad esempio \(E((p\land q) U r)\equiv E((\neg(\neg p \lor \neg q))U r)\). Andremo a etichettare ogni stato che soddisfa \(p\) con \(f_p\), ogni stato che soddisfa \(q\) con \(f_q\), ogni stato che soddisfa \(\neg p \lor \neg q\) con \(f_{\neg p \lor \neg q}\) e così via.

Prima di vedere l'algoritmo però ci serve ampliare il concetto di one-step unfolding visto in precedenza.
In particolare, \(\phi_1 U \phi_2 \equiv \phi_2 \lor (\phi_1 \land X(\phi_1 U \phi_2))\). In \(CTL^*\) quindi avremo
\[E(\phi_1 U \phi_2 )\equiv E(\phi_2 \lor (\phi_1 \land X(\phi_1 U \phi_2))) \equiv E\phi_2 \lor E(\phi_1 \land X(\phi_1 U \phi_2))\]

Nella sintassi di CTL\@, non possiamo esprimere \(E \phi_2\), ma quello che sta dicendo questa formula è che esista un percorso il cui primo stato soddisfa \(\phi_2\), che è equivalente a richiedere semplicemente \(\phi_2\).
Per quanto riguarda \(E(\phi_1 \land X(\phi_1 U \phi_2))\) stiamo richiedendo che esista un percorso in cui \(\phi_1\) è vero in radice e che lo stato successivo soddisfa \(\phi_1 U \phi_2\). Questo è equivalente a \(\phi_1 \land E(X (\phi_1 U \phi_2))\).
Quest'ultimo pezzo può essere riscritto come \(\phi_1 \land EX (E \phi_1 U \phi_2)\), e dunque abbiamo una formula in CTL per one-step unfolding.

Dobbiamo fare lo stesso con \(G\) e per farlo vedere passeremo per \(R\)
\[E(\phi_1 R \phi_2) \equiv E(\phi_2 \land (\phi_1 \lor X(\phi_1 R \phi_2)))\equiv\]
\[\phi_2 \land E(\phi_1 \lor X(\phi_1 R \phi_2))\equiv\]
\[\phi_2 \land (\phi_1 \lor EX (E \phi_1 R \phi_2))\]

Essendo che \(G \phi_2 \equiv \bot R \phi_2\), abbiamo che \(EG \phi_2 \equiv \phi_2 \land EX (EG \phi_2)\) sostituendo e semplificando \(\bot\).

Iniziamo quindi a vedere questi algoritmi partendo dal \(X\).

Iniziamo dalla definizione di chiusura, similmente a quella che abbiamo visto per LTL.

\(cl(\phi) \in NCTL\) è definita come
\begin{itemize}
  \item \(\phi \in cl(\phi)\)
  \item \(\neg \phi' \in cl(\phi)\implies \phi' \in cl(\phi)\)
  \item \(\phi_1 \lor \phi_2 \in cl(\phi)\implies \phi_1, \phi_2 \in cl(\phi)\)
  \item \(EX \phi' \in cl(\phi)\implies \phi' \in cl(\phi)\)
  \item \(E\phi_1 U \phi_2 \in cl(\phi)\implies \phi_1, \phi_2 \in cl(\phi)\)
  \item \(EG \phi' \in cl(\phi)\implies \phi' \in cl(\phi)\)
\end{itemize}

Quello che farà l'algoritmo è etichettare ogni stato con le formule di \(cl(\phi)\) che sono soddisfatte in quello stato andando bottom-up\@, ovvero partendo dalle formule proposizionali e risalendo fino alla formula \(\phi\) che vogliamo verificare.

Per \(EX\) l'algoritmo è il seguente

\begin{algorithmdesc}{Checking \(EX\)}
  \begin{algorithmic}[1]
    \Statex{} \textbf{signature} \(\textsc{CheckEX:} \quad Kripke \times Formula \to Boolean\)
    \Function{\(\textsc{Resolution}\)}{$(k,\phi)$} % chktex 46
    \State{} \(T \gets \set{w \in W}[k,w\models \phi]\)
    \While{\(T\neq \emptyset\)}
    \State{} \(s \gets\) choose \(T\)
    \For{\(t \in R^{-1}(s)\)}
    \State{} \(r(t)\gets r(t)\cup{EX p}\)
    \EndFor{}
    \EndWhile{}
    \EndFunction{}
  \end{algorithmic}
\end{algorithmdesc}

\chapter{Lezione 29/04}

ALTRI ESEMPI SU SPIN\@
